\documentclass[10pt,a4paper]{article}

\usepackage[T1]{fontenc}
\usepackage{amsmath}
\usepackage{amssymb}
\usepackage{amsthm}
\usepackage{array}
\usepackage{booktabs}
\usepackage{geometry}
\usepackage{listings}
\usepackage{longtable}
\usepackage{microtype}
\usepackage{url}
\usepackage[hidelinks]{hyperref}

% Snapshot register: the single place where the audited algorithm
% snapshot is identified. main.tex inputs this file, and the Makefile
% parses it, so keep each definition on its own line in exactly the
% form \newcommand{\Name}{value}.
\newcommand{\SnapshotRelease}{v0.4.0}
\newcommand{\SnapshotCommit}{2a4ef58e3d2e18b01ae67dd1ed4af207e3a79130}
\newcommand{\SnapshotHeaderSha}{ede45489d8e4dc523a97b4bb1565c8b2167d4caf84902ebb35c5645b622208c8}
\newcommand{\SnapshotVerification}{0xF3C4A9B4}
\newcommand{\SnapshotDate}{1 August 2026}


\geometry{margin=27mm}
\setlength{\parindent}{0pt}
\setlength{\parskip}{0.55em}
\setlength{\emergencystretch}{2em}
\lstset{
  basicstyle=\ttfamily\small,
  columns=fullflexible,
  keepspaces=true,
  showstringspaces=false,
  aboveskip=0.6em,
  belowskip=0.6em
}

\newtheorem{lemma}{Lemma}
\newtheorem{theorem}{Theorem}
\newtheorem{proposition}{Proposition}
\newtheorem{corollary}{Corollary}

\newcommand{\rotl}{\operatorname{rotl}}
\newcommand{\load}{\operatorname{load64le}}
\newcommand{\inj}{\operatorname{inj}}
\newcommand{\injb}{\operatorname{inj2}}
\newcolumntype{L}[1]{>{\raggedright\arraybackslash}p{#1}}

\title{Hayahash: A Portable Scalar Non-Cryptographic Hash\\
with 64- and 128-bit Digests}
\author{Ville Vesilehto}
\date{Working draft, \SnapshotDate}

\begin{document}

\maketitle

\begin{abstract}
Hayahash is an experimental, seeded, non-cryptographic hash function for
targets that have ordinary 64-bit multiplication but no cheap 128-bit
product and no guaranteed SIMD.  Its bulk loop carries only an XOR and a
multiplication per lane, over eight independent lanes.  Adjacent input
words are linked by a chained add-and-rotate transform that is bijective
on the whole stripe sequence, so no structural difference pattern can
cancel inside the absorb.  The input length is absorbed in the finalizer,
which makes the digest a pure function of the seed and the bytes read so
far and permits a streaming interface that reproduces the one-shot digest
exactly.  A 128-bit variant extracts a second word from the same state
walk; its low word equals the 64-bit digest for every input and seed, and
on inputs of at most 16 bytes it is injective for each fixed length and
seed.  This paper specifies the current development snapshot, proves these
structural properties, and records the evidence status.  Both digest
widths pass SMHasher3's default suite.  There is no proof for the complete
construction and no cryptographic security claim.
\end{abstract}

\section*{Preface}

This is a living paper for a working algorithm.  The normative artifact is
the C header \cite{vesilehto2026hayahash} as released in
\texttt{\SnapshotRelease}.  The commit below is the one that last changed
the header, which a release commit cannot name for itself, together with
the header's SHA-256 digest:

\begin{center}
\texttt{\SnapshotCommit}\\
\texttt{\SnapshotHeaderSha}
\end{center}

The digest, not the commit, identifies the audited artifact: a commit that
leaves the header byte-identical carries every result in this paper
forward unchanged.  These values are defined once, in
\path{paper/snapshot.tex}; an algorithm change updates the register, the
evaluation records, and this text as one reviewable change, and
\texttt{make -C paper check} verifies that the register, the checkout, and
the header agree.

The \texttt{v0.5} release changed every digest value: the length term
moved from the initial-state premix to the finalizer, the change that
makes the streaming interface possible.  Records produced for earlier
digests describe a superseded function and are never cited as evidence
about the current one.  Every material claim is tied to an entry in the
claim register \path{paper/AUDIT.md}, which states whether it is derived
from the code, reproduced, archived with raw output, reported without a
local reproduction, or open.  Placeholder sections state what must exist
before they can be filled in.

\section{Introduction}

\subsection{Motivation}
\label{sec:motivation}

The fastest current 64-bit hashes obtain their speed from one of two
hardware assumptions: a cheap \(64\times64\to128\)-bit multiplication,
folded into the state by designs in the wyhash line
\cite{wang2021wyhash,nicoshev2025rapidhash}, or SIMD widths negotiated per
target, as in the fast paths of the xxHash family
\cite{collet2023xxhashspec}.  Both assumptions fail on real targets:
baseline WebAssembly has \texttt{i64.mul} but neither a 128-bit product
nor guaranteed SIMD, and portable-C fallbacks, small in-order cores, and
older 64-bit processors are in the same position.  The cost is not
marginal; in the archived WebAssembly record, rapidhash sustains
6.4~GB/s where Hayahash sustains 23.6 in the same runtime on the same
machine (Section~\ref{sec:performance}).

A second constraint is independent of hardware.  A hash deployed across
many language runtimes must produce one exact digest everywhere, so every
required primitive must exist, with identical semantics, in every port.
Unsigned 64-bit addition, XOR, rotation, shifts, and the low half of
multiplication meet that bar; widening multiplies, carry flags, and vector
permutes do not.

Hayahash explores the resulting question: how fast can a hash be when
restricted to portable scalar operations and the low half of 64-bit
multiplication, with alignment-safe reads and a result independent of
host byte order?

\subsection{The construction and why it is shaped this way}
\label{sec:case}

Four decisions carry the design.  Section~\ref{sec:notes} adds the
supporting arguments and Section~\ref{sec:properties} proves what can be
proved.

\textbf{The bulk loop is latency-bound on nothing longer than add and
multiply.}  The throughput limit of a scalar hash loop on a wide core is
the loop-carried chain per lane divided by the number of independent
lanes.  Hayahash runs eight lanes over 64-byte blocks; the carried path
per lane is one XOR and one multiplication, roughly four cycles, with no
cross-lane operation on it.  ChibiHash v2, the closest published function
in the same primitive class, carries roughly five cycles per 8-byte
stripe across four lanes \cite{nrk2024chibihash}.  The archived records
show the consequence: at 1~MiB, 30.7 versus 19.0~GB/s on an Apple M1 Pro
and 61.5 versus 31.1~GB/s on a Zen~5 core.

\textbf{The absorb transform is chained, and the chain is bijective.}
Each stripe is absorbed as \(t_i = w_i + \rotl(w_{i-1},27)\), where
\(w_{i-1}\) is the previous raw input word, carried across lane, block,
and loop boundaries.  The map from word sequence to absorbed sequence is a
bijection (Lemma~\ref{lem:chain}), so no difference pattern can cancel
inside the absorb values; an earlier variant that combined staggered word
copies by XOR was GF(2)-linear and had a nullspace, which produced
constructed collisions.  The rotated copy also plants every stripe's bits
at low positions of the next absorb, where addition and rotation mix two
incompatible algebras: cancelling both copies of a difference requires
matching carries by luck, never by structure.  The rotation applies to a
word already in a register, off the carried path, so the mechanism is
latency-free.

\textbf{The length is absorbed in the finalizer, not the initial state.}
All lane initial values derive from the seed alone, so the state after
\(k\) bytes depends only on the seed and those bytes.  This purity is what
makes a streaming interface with one-shot-identical digests possible; the
previous spelling premixed the length into the initial state, and a
streaming state would have needed the total before the first byte.  The
length enters as \(nK\) inside a finalizer multiplication:
\(n \mapsto nK\) is injective, and Lemma~\ref{lem:length} shows that
inputs of different lengths reaching the finalizer with identical state
still produce different digests in both output words.  That property is
what makes the overlapping tail reads safe; all-zero inputs of lengths 17
through 31 produce identical absorb sequences and are separated by the
length term alone.  A simpler spelling that XORs \(nK\) after the final
multiplications fails SMHasher3's SeedZeroes differentials; the chosen
spelling costs no additional multiplication.

\textbf{Tails read overlapping whole words; short inputs take a dedicated
two-multiply path.}  Inputs longer than 16 bytes are consumed only as
whole 64-bit words, the final reads taken from the end of the input at
overlapping offsets, a technique adopted from wyhash
\cite{wang2021wyhash}; no byte-at-a-time loop exists for any length.
Inputs of at most 16 bytes load two words, spread each through a different
bijective injection, multiply each by a different odd constant, and merge
through a strong finalizer.  For the 128-bit digest this path is injective
for every fixed length and seed (Theorem~\ref{thm:short128}); on exactly
16-byte inputs it is a permutation of the 128-bit space.

Two properties fall out of the same state walk rather than being bolted
on: the 128-bit digest extracts a second word at the end, with its low
word equal to the 64-bit digest and its cost confined to the finalizer,
and the streaming interface reproduces both widths bit for bit under
every split of the input (Proposition~\ref{prop:streaming}).  The design
was hardened by failure: SMHasher3 \cite{wojcik2022smhasher3} found five
structural collision classes in earlier iterations, and every current
anti-collision mechanism is documented with the failure that produced it
(Table~\ref{tab:failures}) and covered by a regression test.

\subsection{Contributions}

\begin{enumerate}
\item An exact specification of both digest widths, the streaming
      interface, and the fixed algorithm parameters
      (Section~\ref{sec:spec}).
\item Proofs of the structural properties the design relies on
      (Section~\ref{sec:properties}).
\item The design arguments behind each mechanism, with the failure classes
      that shaped them (Section~\ref{sec:notes}).
\item An evidence register separating code-derived, reproduced, archived,
      reported, and open claims (Section~\ref{sec:evaluation}).
\end{enumerate}

\subsection{Scope and non-goals}

Hayahash maps a byte string and a 64-bit seed to a 64-bit or 128-bit
result, for hash tables, in-process checksums, and similar non-adversarial
uses.  It is not for passwords, signatures, message authentication,
untrusted hash-table keys, or any use requiring collision or preimage
resistance; the seed is an input, not a security mechanism.  The algorithm
is experimental: constants, thresholds, and digests may change, and
nothing should persist Hayahash values yet.

The design constraints are: unsigned 64-bit scalar arithmetic;
multiplication modulo \(2^{64}\) without the high product; no required
SIMD or architecture-specific extension; alignment-safe reads independent
of host byte order; and one reference mapping across all implementations.
The C source contains per-target scheduling choices; they use the same
operations and are specified to preserve the digest.

\section{Related work}
\label{sec:related}

MurmurHash3 established multiply-and-rotate scalar mixing with a strong
bijective finalizer \cite{appleby2011murmurhash3}; Hayahash reuses its
finalizer constants in the 128-bit output path, and the short-input
finalizer is Evensen's \texttt{moremur} refinement of the same shape
\cite{evensen2018moremur}.  wyhash and rapidhash define the speed frontier
where a widening multiplication is cheap
\cite{wang2021wyhash,nicoshev2025rapidhash}; the xxHash family documents
the scalar and vectorized state of the art \cite{collet2023xxhashspec}.
ChibiHash v2 is the baseline used throughout because it occupies the same
narrow class, a small portable scalar function on the low half of
multiplication \cite{nrk2024chibihash}, so comparisons are between design
choices rather than primitive sets; Hayahash also borrows its 1-to-3-byte
load pattern.

Empirical suites and mathematical guarantees are separate evidence.
SMHasher3 tests distribution and collisions on generated key sets
\cite{wojcik2022smhasher3}; a clean run does not prove collision
resistance, and an earlier Hayahash digest passed the suite while a
constructed collision existed.  PolymurHash states an almost-universality
bound under an explicit model \cite{peters2023polymurhash}; Hayahash has
no comparable bound, and producing one, or documenting why the
construction admits none, is open work.

\section{Notation and interface contract}
\label{sec:contract}

All state values are unsigned 64-bit words.  Addition, multiplication, and
left shift discard bits above bit 63; right shift is logical.
\(\rotl(x,r)\) rotates left by \(r\).  \(\load(B,p)\) reads eight bytes at
offset \(p\), first byte least significant; a four-byte load is
zero-extended.  The implementation realizes loads with \texttt{memcpy}
plus byte swaps on known big-endian hosts, or explicit byte assembly, so
no alignment or endianness assumption reaches the digest.

The constants, all odd, are:

\begin{align*}
K   &= \mathtt{9E3779B97F4A7C15}_{16}, &&\text{\(\lfloor 2^{64}/\varphi
\rfloor\),}\\
M_1 &= \mathtt{3C79AC492BA7B653}_{16}, &&\text{\texttt{moremur},}\\
M_2 &= \mathtt{1C69B3F74AC4AE35}_{16}, &&\text{\texttt{moremur},}\\
N_1 &= \mathtt{FF51AFD7ED558CCD}_{16}, &&\text{MurmurHash3 finalizer,}\\
N_2 &= \mathtt{C4CEB9FE1A85EC53}_{16}, &&\text{MurmurHash3 finalizer.}
\end{align*}

The input length \(n\) must be non-negative with \(n\) readable bytes; a
null pointer is permitted only when \(n=0\).  Two linear word injections
are used:

\begin{align*}
\inj(w)  &= w \oplus \rotl(w,21) \oplus \rotl(w,41),\\
\injb(w) &= w \oplus \rotl(w,11) \oplus \rotl(w,50).
\end{align*}

Three finalizers are used: \texttt{fmix} is the \texttt{moremur} map
\cite{evensen2018moremur}, \texttt{fmix128} is the MurmurHash3 finalizer
\cite{appleby2011murmurhash3}, and \texttt{long\_fmix} is a single-multiply
avalanche:

\begin{lstlisting}
fmix(x):                fmix128(x):             long_fmix(x):
    x ^= x >> 27            x ^= x >> 30            x ^= x >> 37
    x *= M1                 x *= N1                 x *= K
    x ^= x >> 33            x ^= x >> 31            x ^= x >> 32
    x *= M2                 x *= N2                 return x
    x ^= x >> 27            x ^= x >> 33
    return x                return x
\end{lstlisting}

\section{Specification}
\label{sec:spec}

Every path begins with
\[
    s = \mathit{seed} \oplus K, \qquad \mathit{lenmix} = nK .
\]
The length appears nowhere else: the initial state depends on the seed
alone, and \(\mathit{lenmix}\) is consumed by the finalizer.  Inputs of at
most 16 bytes take the short path; 17 to 319 bytes run four lanes over
32-byte blocks; 320 bytes and up run eight lanes over 64-byte blocks, then
fall through to the four-lane machinery for the remainder.

\subsection{Inputs of at most 16 bytes}
\label{sec:spec-short}

Two words \(a\) and \(b\) are formed without reading outside the input:

\begin{lstlisting}
if n >= 8:
    a = load64le(B, 0)
    b = load64le(B, n - 8)
else if n >= 4:
    a = load32le(B, 0)
    b = load32le(B, n - 4)
else if n > 0:
    a = B[0]
    b = (B[n / 2] << 8) | (B[n - 1] << 16)
else:
    a = 0
    b = 0
\end{lstlisting}

The head and end reads overlap for some lengths; Lemma~\ref{lem:length} is
what makes that safe.  The words are spread by different injections,
combined with different seed derivations, and multiplied by different odd
constants:

\begin{align*}
x &= (\inj(a) \oplus s \oplus K)\,K,\\
y &= (\injb(b) \oplus \rotl(s,23) \oplus (K \mathbin{\gg} 19))\,M_1.
\end{align*}

The 64-bit result is
\(\operatorname{fmix}(\rotl(x,27) \oplus y \oplus \mathit{lenmix})\).

\subsection{The absorb chain}
\label{sec:spec-absorb}

Longer inputs are consumed as a sequence of 64-bit words
\(w_0, w_1, \ldots\) called stripes.  With \(w_p\) the previous stripe,
initially zero, one stripe is absorbed into a lane \(h\) as

\begin{lstlisting}
absorb(h, w):
    h = (h ^ (w + rotl(wp, 27))) * K
    wp = w
\end{lstlisting}

The addition happens before the XOR with the lane.  The previous stripe
carries across lane boundaries, block boundaries, and the boundary between
the eight-lane and four-lane loops.

\subsection{Inputs of 17 to 319 bytes}
\label{sec:spec-mid}

The four-lane state derives from \(s\) and shifted copies of \(K\):

\begin{align*}
h_0 &= s \oplus K, &
h_1 &= \rotl(s,17) + (K \mathbin{\ll} 21),\\
h_2 &= \rotl(s,34) \oplus (K \mathbin{\gg} 13), &
h_3 &= \rotl(s,51) + (K \mathbin{\ll} 42).
\end{align*}

Each complete 32-byte block absorbs four consecutive stripes into
\(h_0,h_1,h_2,h_3\) in that order.  The remainder \(r < 32\) goes to the
wall and tail.

\subsection{Inputs of at least 320 bytes}
\label{sec:spec-bulk}

Four more lanes are initialized:

\begin{align*}
h_4 &= s + (K \mathbin{\gg} 27), &
h_5 &= \rotl(s,13) \oplus (K \mathbin{\ll} 9),\\
h_6 &= \rotl(s,26) + (K \mathbin{\gg} 40), &
h_7 &= \rotl(s,39) \oplus (K \mathbin{\ll} 30).
\end{align*}

Each complete 64-byte block absorbs eight consecutive stripes into
\(h_0,\ldots,h_7\); after the eighth absorb, the raw eighth stripe is
added to \(h_0\):
\[
    h_0 \mathrel{+}= w_p .
\]
This checkpoint is part of the digest.  After the last complete block, the
upper lanes fold into the lower:

\begin{align*}
h_0 &\leftarrow (h_0 \oplus \rotl(h_4,11))\,K, &
h_1 &\leftarrow (h_1 \oplus \rotl(h_5,19))\,K,\\
h_2 &\leftarrow (h_2 \oplus \rotl(h_6,31))\,K, &
h_3 &\leftarrow (h_3 \oplus \rotl(h_7,47))\,K.
\end{align*}

If at least 32 bytes remain, one four-lane block runs; the previous-stripe
chain continues across the fold.  The 320-byte threshold is an algorithm
parameter, not a tuning knob: an input of exactly 320 bytes takes the
eight-lane path, moving the threshold changes digests, and its value keeps
the four-lane path shorter than the absorb rotation's 64-stripe orbit
(Section~\ref{sec:notes}).

\subsection{Wall, tail, and finalizer}
\label{sec:spec-tail}

After the complete-block loops, the final stripe's dangling rotated copy
is absorbed; this is called the wall, and it is a no-op when no complete
stripe was processed:
\[
    h_0 \mathrel{+}= \rotl(w_p,27).
\]

With \(p\) the offset of the first unconsumed byte and \(r < 32\) the
remaining length, the tail absorbs whole words through \(\inj\):

\begin{lstlisting}
if r > 16:
    h0 = (h0 + inj(load64le(B, p + 0))) * K
    h1 = (h1 + inj(load64le(B, p + 8))) * K

if r > 0:
    h2 = (h2 + inj(load64le(B, n - 16))) * K
    h3 = (h3 + inj(load64le(B, n - 8))) * K
\end{lstlisting}

The last two reads always cover the final 16 bytes of the input and may
overlap bytes already absorbed; at \(r = 0\) there are no tail reads.  The
finalizer folds the lanes, absorbs the length inside the first
multiplication, and avalanches:

\begin{align*}
t_0 &= (h_0 \oplus \rotl(h_1,13) \oplus \mathit{lenmix})\,K,\\
t_1 &= (h_2 \oplus \rotl(h_3,33))\,K,\\
z   &= s \oplus t_0 \oplus \rotl(t_1,29).
\end{align*}

The 64-bit result is \(\operatorname{long\_fmix}(z)\).

\subsection{The 128-bit digest}
\label{sec:spec-128}

hayahash128 runs the identical state walk and extracts two words at the
end.  On the short path, with
\(u = \rotl(x,27) \oplus y \oplus \mathit{lenmix}\):

\begin{align*}
\mathit{lo} &= \operatorname{fmix}(u),\\
\mathit{hi} &= \operatorname{fmix128}(x + \rotl(u,32)).
\end{align*}

On the longer paths, with \(t_0\) and \(t_1\) as above:

\begin{align*}
\mathit{lo} &= \operatorname{long\_fmix}(s \oplus t_0 \oplus
\rotl(t_1,29)),\\
\mathit{hi} &= \operatorname{fmix128}(\rotl(s,32) \oplus (t_1 +
\rotl(t_0,47))).
\end{align*}

The low word is exactly hayahash64 for every input and seed.  The high
word combines \(t_0\) and \(t_1\) with addition where the low word uses
XOR, so one GF(2)-linear cancellation cannot erase a difference from both
words, and it is finalized by a different bijection.

\subsection{The streaming interface}
\label{sec:spec-streaming}

The streaming state holds the eight lanes, the previous stripe, the seed,
the running total, and a 448-byte buffer with a 128-byte retention floor.
\texttt{update} appends bytes; \texttt{digest} may be called at any time,
does not modify the state, and returns either width.  Bytes are consumed
only in whole 64-byte blocks, in one-shot block order from stream offset
zero; at digest time the buffer retains at least 128 trailing bytes, which
contains every address the tail reads; totals below 448 bytes remain fully
buffered and dispatch to the one-shot function.
Proposition~\ref{prop:streaming} states the resulting guarantee.

\section{Structural properties}
\label{sec:properties}

This section proves the properties the construction relies on.  Each
statement is deliberately narrow, and the section ends by stating what is
not proved.  Arithmetic is modulo \(2^{64}\);
\(\mathbb{W} = \{0,1\}^{64}\).

\begin{lemma}[Absorb chain bijectivity]
\label{lem:chain}
Fix \(w_{-1} \in \mathbb{W}\) and \(m \geq 0\), and define
\(\Phi(w_0,\ldots,w_m) = (t_0,\ldots,t_m)\) with
\(t_i = w_i + \rotl(w_{i-1},27)\).  Then \(\Phi\) is a bijection on
\(\mathbb{W}^{m+1}\), and if two stripe sequences first differ at index
\(j\), their images agree before index \(j\) and differ at index \(j\).
\end{lemma}

\begin{proof}
Given \((t_0,\ldots,t_m)\), recover \(w_0 = t_0 - \rotl(w_{-1},27)\) and
\(w_i = t_i - \rotl(w_{i-1},27)\) in order, so \(\Phi\) is invertible.  If
the sequences agree before index \(j\), the images agree there and the
values \(w_{j-1}\) coincide, so \(t_j\) differs exactly when \(w_j\)
differs, because adding a fixed word is injective.
\end{proof}

Two distinct inputs that disagree in some stripe therefore feed the lane
machinery different absorb values at the first differing stripe.  The
lemma says nothing about what the lane recurrence and folds do afterward.

\begin{lemma}[Injection bijectivity]
\label{lem:inj}
For \(0 < a < b < 64\), the map
\(\iota(w) = w \oplus \rotl(w,a) \oplus \rotl(w,b)\) is a bijection on
\(\mathbb{W}\); in particular \(\inj\) and \(\injb\) are.
\end{lemma}

\begin{proof}
\(\iota\) is GF(2)-linear and commutes with rotation, so it lies in
\(\mathbb{F}_2[z]/(z^{64}+1)\), with \(\iota\) corresponding to
\(p(z) = 1 + z^a + z^b\).  Over \(\mathbb{F}_2\),
\(z^{64}+1 = (z+1)^{64}\), so \(p\) is invertible if and only if
\(p(1) \neq 0\); with three terms, \(p(1) = 1\).
\end{proof}

\begin{lemma}[Finalizer bijectivity]
\label{lem:fmix}
The maps \(x \mapsto x \oplus (x \mathbin{\gg} k)\), \(1 \leq k \leq 63\),
and \(x \mapsto cx\) for odd \(c\) are bijections on \(\mathbb{W}\); hence
\(\operatorname{fmix}\), \(\operatorname{fmix128}\), and
\(\operatorname{long\_fmix}\) are bijections.
\end{lemma}

\begin{proof}
For the shift map the top \(k\) bits of the image equal those of the
argument, and each lower bit is the argument bit XOR an already-recovered
higher bit, so the argument is recovered from the top down.  An odd \(c\)
is a unit modulo \(2^{64}\).  The finalizers are compositions of these
maps.
\end{proof}

\begin{lemma}[Length separation]
\label{lem:length}
Fix a seed, and let two inputs have lengths \(n \neq n'\) with
\(0 \leq n, n' < 2^{63}\).
\begin{enumerate}
\item If both lengths are at most 16 and the inputs load equal word pairs
      \((a,b)\), the 64-bit digests differ, and the 128-bit digests differ
      in the low word.
\item If both lengths are at least 17 and the inputs reach the finalizer
      with equal lane values \((h_0,h_1,h_2,h_3)\), the 64-bit digests
      differ, and the 128-bit digests differ in both words.
\end{enumerate}
\end{lemma}

\begin{proof}
\(K\) is odd, so \(n \mapsto nK\) is injective, and lengths below
\(2^{63}\) are distinct as words: \(\mathit{lenmix} \neq
\mathit{lenmix}'\).

Case 1.  Equal \((a,b)\) and seed give equal \((x,y)\), so the inputs
\(u = \rotl(x,27) \oplus y \oplus \mathit{lenmix}\) and \(u'\) differ, and
\(\operatorname{fmix}\) is a bijection (Lemma~\ref{lem:fmix}).

Case 2.  Equal lane values give equal \(h_0 \oplus \rotl(h_1,13)\), so
\(t_0\) and \(t_0'\) differ: their pre-multiplication values differ and
multiplication by odd \(K\) is injective.  \(t_1 = t_1'\).  The low-word
input \(s \oplus t_0 \oplus \rotl(t_1,29)\) differs and
\(\operatorname{long\_fmix}\) is a bijection.  The high-word input
\(\rotl(s,32) \oplus (t_1 + \rotl(t_0,47))\) differs because rotation and
addition of the fixed \(t_1\) are injective, and \(\operatorname{fmix128}\)
is a bijection.
\end{proof}

The lemma covers exactly the collision class the overlapping reads
create: distinct-length inputs can produce identical loaded words and
absorb sequences, all-zero inputs of lengths 17 through 31 being the
extreme case, and within each path family such pairs never collide.

\begin{theorem}[Short-input 128-bit injectivity]
\label{thm:short128}
Fix a seed and a length \(n \leq 16\).  The map from \(n\)-byte inputs to
\((\mathit{lo},\mathit{hi})\) is injective.
\end{theorem}

\begin{proof}
The map factors as \(B \mapsto (a,b) \mapsto (x,y) \mapsto (u,x) \mapsto
(\mathit{lo},\mathit{hi})\).

\((a,b)\) determines \(B\): for \(8 \leq n \leq 16\), \(a\) holds bytes
\(0..7\) and \(b\) holds bytes \(n-8..n-1\), and \(n-8 \leq 8\), so every
input byte occupies a fixed position of \(a\) or \(b\); for
\(4 \leq n \leq 7\) the same holds with four-byte words; for
\(1 \leq n \leq 3\), \(a = B[0]\) and \(b\) holds \(B[n/2]\) and
\(B[n-1]\) in disjoint bit ranges, and positions \(0\), \(n/2\), \(n-1\)
cover every index below \(n\).  Distinct equal-length inputs therefore
differ in a fixed bit of \((a,b)\).

\((a,b) \mapsto (x,y)\) is a bijection for fixed seed: each coordinate
composes a bijective injection (Lemma~\ref{lem:inj}), an XOR with a
constant, and an odd multiplication (Lemma~\ref{lem:fmix}).
\((x,y) \mapsto (u,x)\) is a bijection for fixed length: \(y\) is
recovered as \(u \oplus \rotl(x,27) \oplus \mathit{lenmix}\).
\((u,x) \mapsto (\operatorname{fmix}(u), \operatorname{fmix128}(x +
\rotl(u,32)))\) is a bijection: recover \(u\) from \(\mathit{lo}\), then
\(x\) from \(\mathit{hi}\).
\end{proof}

\begin{corollary}
For each seed, hayahash128 restricted to inputs of exactly 16 bytes is a
permutation of the 128-bit space.
\end{corollary}

No analogous claim holds for the 64-bit digest alone, nor across distinct
lengths, nor above 16 bytes; a 64-bit digest of arbitrarily long inputs
necessarily has collisions.

\begin{proposition}[Streaming equivalence]
\label{prop:streaming}
For every seed, byte sequence, and partition of the sequence into update
calls, \texttt{hayahash64\_digest} equals \texttt{hayahash64} of the whole
sequence, and \texttt{hayahash128\_digest} equals \texttt{hayahash128}.
\end{proposition}

\begin{proof}[Proof sketch]
Totals below the 448-byte capacity stay contiguous in the buffer and
dispatch to the one-shot function directly.  Otherwise the total exceeds
320 bytes, so the one-shot function would take the eight-lane path, and
the state machine has committed to the same path.  Updates consume input
only in complete 64-byte blocks, in stream order from offset zero, through
the identical absorb-and-checkpoint transform, matching the one-shot
path, which consumes complete blocks while at least 64 bytes remain.  At
digest time the buffered remainder replays the one-shot epilogue on a copy
of the lanes: leftover blocks, fold, at most one 32-byte round, wall,
tail, and finalizer with \(\mathit{lenmix}\) from the running total.  The
128-byte retention floor guarantees the buffered suffix contains the last
16 bytes of the stream, the deepest address the tail reads back.  The
operation sequences are identical and every read returns the same byte.
\end{proof}

A differential test checks this mechanically: every length through 512
under three seeds and five split patterns, plus 1{,}000 randomized cases
through 20~KiB, requiring one-shot and streaming equality for both widths
and \(\mathtt{hayahash128.lo} = \mathtt{hayahash64}\).

\textbf{What is not proved.}  These lemmas are local.  They do not bound
collisions of the complete construction: the lane recurrence, the folds,
and the 64-bit truncation are outside their scope, and Hayahash has no
almost-universality bound and no adversarial guarantee.  Their value is
negative-space: they close the structural collision channels that broke
earlier iterations and delimit exactly where only empirical evidence
exists.

\section{Design notes}
\label{sec:notes}

The header commentary and \path{docs/design.md} carry the full reasoning
and its measurement history; this section records the arguments the
proofs do not cover.

\textbf{Absorb algebra.}  An odd multiplier spreads a difference only
upward, and the top of the word is a fixed point: \(K \equiv 1 \pmod 4\),
so a difference confined to bits 62 and 63 passes through multiplication
by \(K\) unchanged.  A difference planted near the top therefore stays
there, and an XOR-linked absorb could cancel it with a matching later
difference.  The chained rotated copy gives every stripe a second life at
low positions of the next absorb, under addition rather than XOR, so
cancelling both copies requires specific carry patterns rather than a
structural identity.

\textbf{Rotation 27, checkpoint, wall.}  A difference ladder flips bits in
stripe \(i+1\) to cancel, with carry luck, the rotated copy from stripe
\(i\); the residual walks 27 positions per stripe, visiting every position
with period 64 since \(\gcd(27,64)=1\).  Termination requires reaching the
multiplication's fixed region at the top in a same-lane stripe.  With
rotation 21, three hops return to the top (\(3 \cdot 21 = 63\)) and an
earlier iteration produced exact collisions; with 27, every same-lane
revisit stays far from the top for all practical lengths.  The checkpoint
(raw eighth stripe into \(h_0\) each block) was added after a 64-stripe
rotation-orbit family produced deterministic 576-byte collisions; the wall
stops ladders marching off the end of the input; the 320-byte threshold
keeps the four-lane path under 64 stripes, so the full orbit exists only
where the checkpoint interrupts it.

\textbf{Length placement.}  Premixing the length into the initial state
destroys the purity streaming needs; XORing \(nK\) after the final
multiplications leaves modular low-bit structure across lengths that
SMHasher3's SeedZeroes differentials detect.  Absorbing \(nK\) inside the
\(t_0\) multiplication satisfies both constraints at no extra
multiplication.

\textbf{Short path.}  Each loaded word is injected before its
multiplication, so a difference confined to the top bits, which would face
only a few bits of carry luck after an upward-spreading multiply, first
receives low copies.  The two words use different injections and
multipliers because an earlier iteration with one injection on both sides
admitted sparse seed differences erasable by aligned key bits on both
sides at once.

\textbf{Folds.}  Folds combine as \((h \oplus \rotl(h',r))\,K\): an XOR of
multiply-spread values cancels only through carry luck, where an additive
fold could cancel exactly against the additive absorb.  The \(t_1\) fold
uses rotation 33 because 37, the additive inverse of 27 modulo 64,
re-aligned a stripe difference with its rotated copy at the fold.

\small
\begin{longtable}{@{}L{0.25\linewidth}L{0.30\linewidth}L{0.34\linewidth}@{}}
\caption{Documented structural failures and current responses.}
\label{tab:failures}
\\
\toprule
Earlier form & Observed problem & Current response \\
\midrule
\endfirsthead
\toprule
Earlier form & Observed problem & Current response \\
\midrule
\endhead
XOR of staggered stripe copies &
GF(2)-linear nullspace across the absorb sequence &
Chained addition of the previous rotated raw stripe
(Lemma~\ref{lem:chain}). \\
\addlinespace
Same injection on both short-path words &
Sparse seed differences erasable by aligned key bits on both sides &
Different injections and multipliers for \(a\) and \(b\). \\
\addlinespace
Absorb rotation 21, no wall &
Carry-matched difference ladders could terminate at the top of the word
or exit past the last stripe &
Rotation 27; wall absorb of the final rotated copy. \\
\addlinespace
Fold rotation 37 &
Re-aligned with the absorb rotation and allowed XOR cancellation at the
fold &
Fold rotation 33. \\
\addlinespace
No bulk checkpoint &
64-stripe rotation-orbit family gave deterministic 576-byte collisions &
Raw eighth stripe added to \(h_0\) after every block; deterministic
regression retained. \\
\bottomrule
\end{longtable}
\normalsize

\section{Evaluation}
\label{sec:evaluation}

The claim register \path{paper/AUDIT.md} classifies each material
statement as code-derived, reproduced, archived, reported, or open; raw
outputs live under \path{paper/results} with provenance headers.

\subsection{Quality}
\label{sec:quality}

The local harness is deterministic (fixed-seeded), so its numbers depend
on the digest, not the host: strict-avalanche measurements over input and
seed bits (4{,}000 trials per cell, warn above 0.03, fail above 0.045) and
exact collision checks over 24 structured key sets that reproduce every
failure class of Table~\ref{tab:failures}.  At this snapshot the worst
input-bit bias is 0.0360 and the worst seed-bit bias 0.0330, neither near
the failure threshold; all 24 sets are clean, and the streaming
differential of Proposition~\ref{prop:streaming} passes in the same run
(archived in \path{paper/results/quality.txt}).  As a control, the same
harness finds 512 collisions in ChibiHash v2's rotation-orbit set: a
constructed-set result, not a quality ranking, showing that statistical
suites and construction-specific tests serve different roles.

\subsection{SMHasher3}
\label{sec:smhasher}

The canonical verification values are reproduced directly from the
reference header by \path{paper/tools/reference_check.c}; the byte-swapped
values are registered in the SMHasher3 translation and exercised by the
tool's self-test in every run:

\begin{center}
\begin{tabular}{@{}lcc@{}}
\toprule
 & canonical (LE) & byte-swapped (BE) \\
\midrule
hayahash64  & \texttt{\SnapshotVerification} & \texttt{0x805DE5C0} \\
hayahash128 & \texttt{\SnapshotVerificationWide} & \texttt{0x46140A64} \\
\bottomrule
\end{tabular}
\end{center}

The harness in \path{tests/smhasher3} pins upstream commit
\texttt{51d3cd1a} and installs a self-contained translation written for
upstream submission; its agreement with the header is enforced by the
verification values and the repository's conformance machinery.  Archived
for the current digest: the full default suite passes 188 of 188 for
hayahash128 on an EPYC 9655 at that pin
(\path{paper/results/epyc9655-smhasher3-128-conformance.txt}).  Reported,
without an archived raw dump: both widths pass the full suite on an Apple
M1 Pro.  The previous digest's nine-build sweep across four hosts and five
compilers demonstrates the verification pipeline but says nothing about
the current digest; repeating it is open work.  A clean run is a
black-box statistical result: it neither proves collision resistance nor
ranks passing functions, as this project's own constructed collision after
a clean run demonstrated.

\subsection{Performance (TBD)}
\label{sec:performance}

This draft deliberately carries no benchmark tables.  The archived records
under \path{paper/results} support the following summary, which is what
the tables must substantiate when they return:

\begin{itemize}
\item hayahash64 leads ChibiHash v2, the closest same-class function, in
      every cell of the two archived bare-metal records (Apple M1 Pro and
      Zen 5): bulk throughput, seed-chained latency, and independent
      throughput, from 1 byte to 1~MiB.  The statement was verified
      mechanically over both records.
\item hayahash128 reaches bulk parity with hayahash64 on every recorded
      host; producing the second word costs under 2.3~ns in the recorded
      small-input cells.
\item On baseline wasm32, without SIMD or a widening multiply, hayahash
      sustains 23.6~GB/s where rapidhash sustains 6.4 in the same runtime,
      consistent with the primitive-class argument of
      Section~\ref{sec:motivation}.
\item Records carry source checksums, host and compiler provenance, and a
      stated method: medians of nine calibrated samples through
      non-inlined wrappers, candidate order rotated.
\end{itemize}

The tables return when: the M1 Pro record is retaken at the released
header (the current one measured a documented output-identical
development state); the harness retains raw samples and an environment
manifest rather than medians alone; and the transcription is generated
from the records instead of copied.

\subsection{Conformance}
\label{sec:conformance}

The C reference is ported to Rust, Go, Zig, Java, C\#, Python, Swift,
JavaScript, and MIPS64 assembly, each covering both widths and asserting
the shared vectors, the verification value, and
\(\mathtt{hayahash128.lo} = \mathtt{hayahash64}\).  CI replays a nightly
C-reference differential corpus through every port and covers big-endian
(s390x under emulation), ILP32 (wasm32), and MSVC x64.  Rerun for this
draft: the reference check for both widths and both verification values,
the quality harness, and the Go, Rust, and Python suites
(\path{paper/results/conformance.txt}).

\section{Limitations and open work}
\label{sec:limitations}

Standing properties: the algorithm is experimental and digests may change;
there is no cryptographic or adversarial claim of any kind; the structural
lemmas do not bound the complete construction; the harness and SMHasher3
are statistical evidence only; the streaming state buffers up to 448
bytes.

Open work, each item filled in only with its artifacts: an explicit model
of the absorb recurrence and folds, and either an almost-universality
bound comparable in form to PolymurHash's \cite{peters2023polymurhash} or
a documented argument why none exists; a rerunnable archived differential
search over the absorb chain; a multi-host SMHasher3 sweep and archived
raw full-suite runs for the current digest; the benchmark completion
criteria of Section~\ref{sec:performance}; and one generated
machine-readable vector table consumed by every port.

\section{Reproducibility}
\label{sec:repro}

From the repository root at the paper snapshot:

\begin{lstlisting}
make -C paper check          # reference vectors, quality, snapshot
make -C tests/smhasher3 run  # full suite; HASH=hayahash128 for the wide one
make -C paper                # build this PDF with Tectonic
\end{lstlisting}

A snapshot change must update the register, the vectors, the quality
output, and the claim register as one reviewable change; the procedure is
in \path{paper/README.md}.

\section{Conclusion}

Hayahash asks how fast a hash can be on portable scalar operations and the
low half of 64-bit multiplication, under one exact digest across every
implementation.  The answer is a bulk loop carrying one XOR and one
multiplication across eight lanes, an absorb chain that is bijective on
the stripe sequence, a length term in the finalizer that permits an exact
streaming interface, and a 128-bit extraction whose low word is the 64-bit
digest.  Every anti-collision mechanism is tied to a documented failure
and a regression.  The evidence justifies continued work, not security or
universal quality claims, and the paper is structured so that stronger
claims can be added only together with their artifacts.

\section*{Revision history}

\begin{tabular}{@{}L{0.16\linewidth}L{0.18\linewidth}L{0.58\linewidth}@{}}
\toprule
Date & Paper snapshot & Change \\
\midrule
2026-08-09 & \texttt{v0.5.0} & Restructured and condensed: preface,
introduction with the design case, exact specification of both widths and
the streaming interface, proved structural properties, compact design
notes.  Snapshot advanced to the \texttt{v0.5} digest with regenerated
vectors and a rerun quality harness; benchmark tables replaced by a
placeholder with completion criteria; references trimmed to load-bearing
entries. \\
\addlinespace
2026-08-01 & \texttt{v0.4.0} & Snapshot advanced to \texttt{v0.4.0} after
the fifth optimization pass rewrote per-compiler dispatch without moving
the digest.  SMHasher3 rerun on builds covering every compiled shape. \\
\addlinespace
2026-07-31 & \texttt{v0.3.0} & SMHasher3 section filled in with a pinned
adapter and archived runs; snapshot register introduced as the single
source of the audited identity. \\
\addlinespace
2026-07-31 & first and second drafts & Exact specification, audit
register, local quality results, archived M1 and Zen 5 comparisons;
restructure into a living document with placeholder sections carrying
completion criteria. \\
\bottomrule
\end{tabular}

\bibliographystyle{plain}
\bibliography{references}

\end{document}
